\section{Modelo de datos del context: qué debe almacenar realmente la socialización con IA}

Para que el avatar se parezca al usuario, una ficha de perfil no basta. La socialización con IA necesita un context de varias capas: identidad estable, preferencias de largo plazo, estado reciente, estilo de expresión, objetivos sociales, grafo de relaciones, límites de privacidad e historial de retroalimentación. Cada capa tiene una vigencia temporal, una sensibilidad y una forma de autorización distintas, por lo que no pueden mezclarse en un único campo de «memoria».

\subsection{Capas del context}

\noindent\begin{longtable}{p{0.22\textwidth}p{0.34\textwidth}p{0.35\textwidth}}
\toprule
Capa & Contenido & Tratamiento en el producto \\
\midrule
Capa de identidad & Identidad básica, profesión, ciudad, profile público & Puede ser pública y editable; sensibilidad baja. \\
Capa de preferencias & Intereses, gusto estético, formas de expresión que agradan o desagradan & Se usa para las recomendaciones y para generar el tono. \\
Capa de relaciones & Amigos, seguidos, interacciones previas, vínculos débiles & Alto valor y también alta sensibilidad; requiere permisos. \\
Capa de estado & Emociones recientes, objetivos, actividades en curso & Muy sujeta al tiempo; requiere un mecanismo de caducidad. \\
Capa de expresión & Estilo de escritura, humor, opiniones, límites & Determina si el avatar se parece o no al usuario. \\
Capa de retroalimentación & Aprobación o rechazo del usuario a las acciones del avatar & Base para entrenar al avatar y ajustar la estrategia. \\
\bottomrule
\end{longtable}

\begin{importantbox}{Con el context no importa tener más, sino tenerlo mejor separado en capas}
Cada capa del context difiere en sensibilidad, vigencia temporal y uso. Si un producto de socialización con IA no las gestiona por separado, enfrentará al mismo tiempo riesgos de privacidad y una experiencia confusa.
\end{importantbox}

\subsection{Ciclo de vida del context}

El context debe tener un ciclo de vida: recopilación, interpretación, uso, retroalimentación, caducidad y eliminación. Por ejemplo, «hoy quiero conocer a gente que hace productos de IA» es un objetivo de corto plazo y no debe fijarse de forma permanente; «me desagrada el humor ofensivo» es una preferencia de largo plazo y puede conservarse por mucho tiempo; en cambio, «el contenido de cierta conversación privada» quizá solo pueda usarse de forma local y no debe entrar en el emparejamiento público.

\begin{warningbox}{El problema de la contaminación de la memoria}
Si el avatar toma una expresión momentánea como una preferencia de largo plazo, o una broma como un valor personal, se produce contaminación de la memoria. Los productos de socialización con IA deben permitir que el usuario corrija y elimine el context; de lo contrario, cuanto más se usen, menos se parecerán a él.
\end{warningbox}

\subsection{Resumen de la sección}

El context es el modelo de datos subyacente de la socialización con IA. Requiere separación en capas, autorización, caducidad y retroalimentación, no una simple acumulación de historiales de chat.
